\section{Genetic information as readout arithmetic: discretization, coarse graining, and correction}
\label{sec:genetic_readout}

HPA--$\Omega$ is basis-independent at the interface level: ``scan--projection readout'' is a structural constraint, not a particular physical substrate.
This motivates a biological interface hypothesis: persistent biological information systems may exhibit architectures that look like finite-resolution readout plus active correction.

\subsection{Interface hypothesis: genetic systems as finite-resolution readout--correction channels}
\label{sec:genetic_hypothesis}

\begin{remark}[Interface hypothesis for genetic coding]
\label{rem:genetic_interface}
The genetic alphabet and its decoding pipeline are modeled as a finite-resolution readout map from a large microscopic configuration space (molecular conformations, binding fluctuations, chemical noise) to a discrete executable output space (amino-acid sequences and regulatory actions).
Degeneracy of the code (many codons mapping to one amino acid) is a coarse-graining that increases robustness under readout noise, while proofreading and repair mechanisms provide active correction to maintain low effective mismatch \cite{Alberts2014,Crick1966,Hopfield1974,Ninio1975,Lindahl1993}.
\end{remark}

We treat this as an interface-level mapping whose scientific content is carried by operational predictions and measurable proxies.
It does, however, align with three basic observations:
\begin{itemize}
  \item \textbf{Finite alphabet with combinatorial capacity.}
  A small discrete alphabet supports large effective state spaces through composition (e.g.\ triplet codons).
  \item \textbf{Degeneracy as controlled coarse graining.}
  Many-to-one mappings (including wobble pairing) trade fine detail for robustness \cite{Crick1966}.
  \item \textbf{Correction is not free.}
  Proofreading and repair require energy and physical transport; therefore their minimal costs must include a Landauer scale and an architecture-dependent geometric impedance (Eq.~\eqref{eq:geom_landauer}).
\end{itemize}

\subsection{Genetic memory as compressed predictive prior}
\label{sec:genetic_prior}

In the AEC language, genetic information is not a static ``blueprint'' but a compressed prior that supports prediction and control across time:
it constrains the space of internal models $M$ that an organism can instantiate and maintain under finite resources.
Maintenance of that prior (replication fidelity, error correction, epigenetic stabilization) consumes dissipation budget $\dot W_{\mathrm{diss}}$ and therefore competes with other uses of energy such as growth and reproduction.
This makes genetics part of the same optimization problem as behavior: allocate limited dissipation and limited predictive information capacity to sustain $\eta_{\mathrm{pred}}$ in the environment.


